\subsection{逆紧映射}

若 \(f: \mathbf{X} \to \mathbf{Y}\) 与 \(f': \mathbf{X}' \to \mathbf{Y}'\) 是两个连续闭映射，则积 \(f \times f': \mathbf{X} \times \mathbf{X}' \to \mathbf{Y} \times \mathbf{Y}'\) 未必是闭映射，即使 \(f\) 形如 \(\iota_{\mathbf{X}}\)。

例. 到豪斯多夫空间中的每个常值映射都是闭的。但若 \(f\) 是常值映射 \(\mathbf{Q} \to 0\)，则 \(f \times \iota_{\mathbf{Q}}\) 是映射 \((x, y) \to (0, y)\)（\(\mathbf{Q}^2\) 到 \(\mathbf{Q}^2\) 中的），即第二个投影，它不是闭的（§ 4，第 2 目，注 1）。

定义 1. 设 \(f\) 是拓扑空间 \(\mathbf{X}\) 到拓扑空间 \(\mathbf{Y}\) 中的一个映射。称 \(f\) 是逆紧的，如果 \(f\) 是连续的，并且映射 \(f \times \iota_{\mathbf{Z}} : \mathbf{X} \times \mathbf{Z} \to \mathbf{Y} \times \mathbf{Z}\) 对每个拓扑空间 \(\mathbf{Z}\) 都是闭的。

我们将在第 2 目与第 3 目中给出逆紧映射的其他刻画。

如果在定义 1 中取空间 Z 由一个点组成，便得到：

命题 1. 每个逆紧映射都是闭的。

命题 2. 设 \(f: \mathbf{X} \to \mathbf{Y}\) 是一个连续单射。则以下三个命题等价：

\begin{enumerate}
\def\labelenumi{\alph{enumi})}
\item
  \(f\) 是逆紧的。
\item
  \(f\) 是闭的。
\item
  \(f\) 是 \(\mathbf{X}\) 到 \(\mathbf{Y}\) 的一个闭子集上的一个同胚。
\end{enumerate}

我们刚才已见 a) 蕴含 b)。由于等价关系 \(f(x) = f(x')\) 就是相等关系，X 关于这个关系的商空间可与 X 等同；故由 § 5，第 2 目，命题 3，b) 蕴含 c)。最后，若 c) 成立，则 \(f \times \iota_{\mathbf{Z}}\) 是 \(\mathbf{X} \times \mathbf{Z}\) 到 \(\mathbf{Y} \times \mathbf{Z}\) 的一个闭子空间上的一个同胚，因而是闭映射；故 c) 蕴含 a)。

命题 3. 设 \(f: \mathbf{X} \to \mathbf{Y}\) 是一个连续映射。若 \(\mathbf{T}\) 是 \(\mathbf{Y}\) 的任意子集，用 \(f_{\mathbf{T}}\) 表示映射 \(\overline{f}^{1}(\mathbf{T}) \to \mathbf{T}\)，它与 \(f\) 在 \(\overline{f}^{1}(\mathbf{T})\) 上取值相同。

\begin{enumerate}
\def\labelenumi{\alph{enumi})}
\item
  若 \(f\) 逆紧，则 \(f_{\mathbf{T}}\) 也逆紧。
\item
  设 \((\mathrm{T}(\iota))_{\iota\in I}\) 是 Y 的一个子集族，其诸内部覆盖 Y，或者它是 Y 的一个局部有限闭覆盖；则若诸映射 \(f_{\mathbf{T}(\iota)}\) 都逆紧，\(f\) 也逆紧。
\end{enumerate}

设 \(\mathbf{Z}\) 是一个拓扑空间。若 \(\mathbf{T}\) 是 \(\mathbf{Y}\) 的任意子集，我们有

\[
f _ {\mathbf {T}} \times \iota_ {\mathbf {Z}} = (f \times \iota_ {\mathbf {Z}}) _ {\mathbf {T} \times \mathbf {Z}};
\]

若 \(f\) 逆紧，则 \(f \times \iota_{\mathbf{Z}}\) 是闭的，因而 \((f \times \iota_{\mathbf{Z}})_{\mathbf{T} \times \mathbf{Z}}\) 也是闭的{[}§ 5，第 1 目，命题 2 a){]}，于是 a) 得证。现在若 \((\mathrm{T}(\iota))_{\iota \in \mathbf{I}}\) 满足 b) 中所述两个条件之一，则覆盖 \((\mathrm{T}(\iota) \times \mathbf{Z})_{\iota \in \mathbf{I}}\)（\(\mathbf{Y} \times \mathbf{Z}\) 的）具有同样的性质；若诸 \(f_{\mathbf{T}(\iota)}\) 逆紧，则诸映射

\[
(f \times \iota_ {\mathbf {Z}}) _ {\mathbf {T} (\iota) \times \mathbf {Z}}
\]

都是闭的，于是 \(f \times \iota_{\mathbf{Z}}\) 是闭的{[}§ 5，第 1 目，命题 2 b){]}。证明完毕。

命题 4. 设 I 是有限集，且对每个 \(i \in I\)，设 \(f_i: X_i \to Y_i\) 是一个连续映射。设 \(X = \prod_{i \in I} X_i\)，\(Y = \prod_{i \in I} Y_i\)，并设 \(f: X \to Y\)

是积映射 \((x_{i})\to (f_{i}(x_{i}))\)。则：

\begin{enumerate}
\def\labelenumi{\alph{enumi})}
\item
  若每个 \(f_i\) 都逆紧，则 \(f\) 逆紧。
\item
  若 \(f\) 逆紧且诸 \(X_i\) 非空，则每个 \(f_i\) 都逆紧。
\end{enumerate}

（我们将在第 2 目，定理 1，推论 3 中看到，这个命题可以推广到无穷积。）

用归纳法，只需考虑 \(\mathbf{I} = \{\mathbf{1},\mathbf{2}\}\) 的情形。a) 设 \(f_{1}, f_{2}\) 逆紧，并设 \(\mathbf{Z}\) 是一个拓扑空间；\(f_{1} \times f_{2} \times \iota_{\mathbf{Z}}\) 是 \(\iota_{\mathbf{Y}_1} \times f_2 \times \iota_{\mathbf{Z}}\) 与

\[
f _ {1} \times \iota_ {\mathbf {X} _ {2}} \times \iota_ {\mathbf {Z}};
\]

的复合；由假设这两个映射都是闭的，故 \(f_{1} \times f_{2} \times \iota_{\mathbf{Z}}\) 也是闭的{[}§ 5，第 1 目，命题 1 a){]}，于是 \(f_{1} \times f_{2}\) 逆紧。

\begin{enumerate}
\def\labelenumi{\alph{enumi})}
\setcounter{enumi}{1}
\tightlist
\item
  现在设 \(f\) 逆紧。设 \(\mathbf{F}\) 是 \(\mathbf{X}_2 \times \mathbf{Z}\) 的一个闭子集，\(\mathbf{G}\) 是 \(\mathbf{F}\) 在 \(\mathbf{Y}_2 \times \mathbf{Z}\) 中（在映射 \(f_2 \times \iota_{\mathbf{Z}}\) 下）的象。则 \(\mathbf{X}_1 \times \mathbf{F}\) 在 \(\mathbf{Y}_1 \times \mathbf{Y}_2 \times \mathbf{Z}\) 中（在 \(f_1 \times f_2 \times \iota_{\mathbf{Z}}\) 下）的象是 \(f_1(\mathbf{X}_1) \times \mathbf{G}\)。由假设，它在 \(\mathbf{Y}_1 \times \mathbf{Y}_2 \times \mathbf{Z}\) 中是闭的；若 \(\mathbf{X}_1 \neq \emptyset\)，则 \(f_1(\mathbf{X}_1)\) 非空，这蕴含 \(\mathbf{G}\) 在 \(\mathbf{Y}_2 \times \mathbf{Z}\) 中是闭的（§ 4，第 3 目，命题 7 的推论）；故 \(f_2\) 逆紧。类似地，\(f_1\) 在 \(\mathbf{X}_2 \neq \emptyset\) 时逆紧。
\end{enumerate}

命题 5. 设 \(f: \mathbf{X} \to \mathbf{X}'\) 与 \(g: \mathbf{X}' \to \mathbf{X}''\) 是两个连续映射。

\begin{enumerate}
\def\labelenumi{\alph{enumi})}
\item
  若 \(f\) 与 \(g\) 都逆紧，则 \(g \circ f\) 逆紧。
\item
  若 \(g \circ f\) 逆紧且 \(f\) 是满射，则 \(g\) 逆紧。
\item
  若 \(g \circ f\) 逆紧且 \(g\) 是单射，则 \(f\) 逆紧。
\item
  若 \(g \circ f\) 逆紧且 \(X'\) 是豪斯多夫的，则 \(f\) 逆紧。
\end{enumerate}

设 \(\mathbf{Z}\) 是一个拓扑空间。我们有

\[
(g \circ f) \times \iota_ {\mathbf {Z}} = (g \times \iota_ {\mathbf {Z}}) \circ (f \times \iota_ {\mathbf {Z}});
\]

若 \(f\) 与 \(g\) 逆紧，则 \(f \times \iota_{\mathbf{Z}}\) 与 \(g \times \iota_{\mathbf{Z}}\) 都是闭的；故{[}§ 5，第 1 目，命题 1 a){]} ( \(g \circ f\) ) × \(\iota_{\mathbf{Z}}\) 是闭的；这就证明了 a)。b){[}相应地，c){]}的证明沿着同样的思路进行，利用 § 5 第 1 目命题 1 的 b){[}相应地，c){]}，并注意到若 \(f\) 是满射（相应地，若 \(g\) 是单射），则 \(f \times \iota_{\mathbf{Z}}\) 是满射（相应地，\(g \times \iota_{\mathbf{Z}}\) 是单射）。最后，为证 d)，考虑交换图

\[
\mathbf {X} \stackrel {\varphi} {\rightarrow} \mathbf {X} \times \mathbf {X} ^ {\prime}\tag{1}
\]

\[
\begin{array}{c} f \downarrow \\ \mathbf {X} ^ {\prime} \xrightarrow [ \psi ]{} \mathbf {X} ^ {\prime \prime} \end{array} \begin{array}{c} \downarrow (g \circ f) \times \iota_ {\mathbf {x} ^ {\prime}} \\ \times \mathbf {X} ^ {\prime} \end{array}
\]

其中 \(\varphi(x) = (x, f(x))\)，\(\psi(x') = (g(x'), x')\)。映射 \(\varphi\)（相应地，\(\psi\)）是 X（相应地，X'）到 \(f\) 的图象（相应地，\(g\) 的图象的镜象）上的一个同胚（\(\S 4\)，第 1 目，命题 1，推论 2）。此外，由于 X' 是豪斯多夫的，图象 \(\varphi(X)\)（即 \(f\) 的图象）在 X × X' 中是闭的（\(\S 8\)，第 1 目，命题 2，推论 2）。故（命题 2）\(\varphi\) 是逆紧的；另一方面，命题 4 表明 ( \(g \circ f\) ) × \(\iota_{X'}\) 是逆紧的。由上面的 a) 及图 (1) 的交换性，\(\psi \circ f\) 是逆紧的；但 \(\psi\) 是单射，故由上面的 c)，\(f\) 是逆紧的。

注. 若 \(\mathbf{X}'\) 不是豪斯多夫的，则可能 \(g \circ f\) 逆紧而 \(f\) 不逆紧；例如，取 \(\mathbf{X}\) 与 \(\mathbf{X}''\) 各由一点组成，\(\mathbf{X}'\) 由两点组成并带最粗拓扑。

推论 1. 若 \(f: \mathbf{X} \to \mathbf{Y}\) 是逆紧映射，则 \(f\) 在闭子集 \(\mathbf{F}\)（\(\mathbf{X}\) 的）上的限制是 \(\mathbf{F}\) 到 \(\mathbf{Y}\) 中的一个逆紧映射。

因为这个限制是复合 \(f \circ j\)，其中 \(j: \mathbf{F} \to \mathbf{X}\) 是典范单射，由命题 2 它是逆紧的。

推论 2. 设 \(f: \mathbf{X} \to \mathbf{Y}\) 是逆紧映射，其中 \(\mathbf{X}\) 是豪斯多夫的。则子空间 \(f(\mathbf{X})\)（\(\mathbf{Y}\) 的）是豪斯多夫的。

由命题 5 c)，只需考虑 \(f(\mathbf{X}) = \mathbf{Y}\) 的情形。此时 \(\mathbf{Y} \times \mathbf{Y}\) 的对角线是 \(f \times f\) 下 \(\mathbf{X}\) 的对角线的象，后者是闭的（§ 8，第 1 目，命题 1）；\(f \times f\) 是逆紧的（命题 4）；故 \(\mathbf{Y} \times \mathbf{Y}\) 的对角线是闭的（命题 1），因此 \(\mathbf{Y}\) 是豪斯多夫的（§ 8，第 1 目，命题 1）。

推论 3. 设 I 是有限集，且对每个 \(i \in I\)，设 \(f_i: X \to Y_i\) 是逆紧映射。若 X 是豪斯多夫的，则映射 \(x \to (f_i(x))\)（X 到 \(\prod_{i \in I} Y_i\) 中的）是逆紧的。

这个映射是积映射 \((x_{i})\to (f_{i}(x_{i}))\)（\(\mathbf{X}^{\mathbf{I}}\) 到 \(\prod_{i}\mathbf{Y}_{i}\) 中的）与 \(\mathbf{X}\) 到 \(\mathbf{X}^{\mathbf{I}}\) 中的对角映射的复合；由于后者是逆紧的（由命题 2 及 § 8，第 1 目，命题 1），结论由命题 4 与命题 5 a) 得出。

推论 4. 设 X 与 Y 是两个拓扑空间，\(f: \mathbf{X} \to \mathbf{Y}\) 是一个连续映射，R 是 X 上的等价关系 \(f(x) = f(y)\)，而

\[
\mathbf {X} \xrightarrow {p} \mathbf {X} / \mathbf {R} \xrightarrow {h} f (\mathbf {X}) \xrightarrow {i} \mathbf {Y}
\]

是 \(f\) 的典范分解。则 \(f\) 逆紧当且仅当 \(p\) 逆紧、\(h\) 是一个同胚且 \(f(\mathbf{X})\) 是 \(\mathbf{Y}\) 的一个闭子集。

由命题 5 a) 与命题 2，这些条件是充分的。反之，若 \(f\) 逆紧，则 \(f\) 是闭的；故 \(f(\mathbf{X})\) 在 \(\mathbf{Y}\) 中是闭的，且 \(h\) 是一个同胚（§ 5，第 2 目，命题 3）；又由命题 5 c)，\(h \circ p\) 是逆紧的；故由命题 5 a)，\(p = \overline{h} \circ (h \circ p)\) 是逆紧的。

